% Subject Instructions Template
% Compile with pdflatex
%
% Guidelines:
% - Maximum 3 pages
% - Pretest with non-experts before running
% - Pretesters should be able to answer:
%   (1) What decision am I making?
%   (2) How do I make it?
%   (3) Where do I record it?
%   (4) What are possible outcomes?

\documentclass[12pt]{article}
\usepackage[left=1in,right=1in,top=1in,bottom=1in]{geometry}
\usepackage{setspace}
\onehalfspacing
\usepackage{enumitem}
\usepackage{booktabs}

\begin{document}

\begin{center}
{\Large \textbf{Instructions}} \\[0.5em]
{\normalsize [Experiment Name]} \\[0.5em]
{\small Please read these instructions carefully. Your payment depends on your decisions.}
\end{center}

\bigskip

\section*{Overview}

% Describe the experiment in plain language.
% What will the participant do? How long will it take?
% What is the purpose (in general terms — no deception, but vagueness is OK)?

[Overview text here.]

\section*{Your Task}

% Describe exactly what decision(s) the participant makes.
% Use concrete language and examples.
% If there are multiple rounds/stages, describe each.

[Task description here.]

\section*{How You Will Be Paid}

% Describe the payment structure clearly.
% - Show-up fee / base payment
% - How decisions translate to earnings
% - Which round(s) count for payment
% - Expected range of total payment

You will receive a base payment of \$[X] for participating. In addition, [describe performance-based payment].

[If random round payment:] At the end of the experiment, one round will be randomly selected to determine your bonus payment. Because any round could be the one that counts, it is in your best interest to make each decision as if it were the one that determines your payment.

\section*{Important Rules}

\begin{itemize}[nosep]
    \item There are no right or wrong answers. We are interested in your genuine preferences.
    \item Your decisions are private. No other participant will know what you chose.
    \item Please do not communicate with other participants during the experiment.
    \item If you have questions at any time, please raise your hand. % [or: use the chat function]
\end{itemize}

\section*{Example}

% Provide a concrete, worked example.
% Walk through a decision and show how payment would be calculated.
% Use different numbers than the actual experiment to avoid anchoring.

[Worked example here.]

\section*{Summary}

\begin{enumerate}[nosep]
    \item You will make [N] decisions.
    \item Each decision involves [brief description].
    \item Your payment depends on [brief description].
    \item The experiment takes approximately [X] minutes.
\end{enumerate}

\bigskip
\begin{center}
{\small If you have any questions, please ask the experimenter now.} \\[0.5em]
{\small When you are ready, please click ``Begin'' to start the comprehension quiz.}
\end{center}

\end{document}
